\section%
%Einheiten mit siunitx-Paket: $\SI{512}{\mega \byte}$, $\SI{235}{\Npsmm}$ \& $l_{ges} = \SI{1234}{\mm}$
%
%Commands für gängige Programme-Namen $\rightarrow$ besonders für Überschriften: \labview und \matlab
%
%Eingekreister Text funktioniert so: \circText{1}, oder so \circTextBig{2}, oder so: \circTextSmall{3}
%
%Ableitungen funktionieren so: $\deriv[2]{s}{t}$, oder so: $\deriv{I}{t}$
%
%Römische Zahlen funktionieren so: \RN{4}
%
%Toleranzangaben funktionieren so: \tol{25}{\mm}{0,2}{0,1}
%
%Versordarstellung mit Steinmetz-Paket: $X_c = \SI{470}{\ohm}\phase{\gradU{-90}}$
%
%Kürzen von Einheiten geht so: $\Delta \varphi=\frac{t}{T}*\gradU{360}=\frac{\SI{-0,00013}{\cancel\s}}{\SI{0,001}{\cancel\s}}*\gradU{360}=\gradU{-46,8}$
%
%Für Tabellen ist die Verwendung des inkludierten booktabs-Paketes zu empfehlen

%\begin{figure}[H]
%	\begin{center}
%		\includegraphics[width=0.85\textwidth]{bilder/bild1.png}
%		\caption{Caption eines Bildes 1}
%		\label{fig:asdf1}
%	\end{center}	
%\end{figure}


%\begin{minipage}{0.5\textwidth}
%	\centering
%	\begin{figure}[H]
%		\begin{center}
%			\includegraphics[height=0.15\textheight]{bilder/bild1.png}
%			\caption{Caption eines Bildes 2}
%			\label{fig:asdf2}
%		\end{center}
%	\end{figure}
%\end{minipage}
%\hfill
%\begin{minipage}{0.5\textwidth}
%	\centering
%	\begin{figure}[H]
%		\begin{center}
%			\includegraphics[height=0.15\textheight]{bilder/bild2.png}
%			\caption{Caption eines Bildes 3}
%			\label{fig:asdf3}
%		\end{center}
%	\end{figure}
%\end{minipage}



%\label{fig:Kapitelnummer-Abbildungsname}	für Abbildungen
%\label{tbl:Kapitelnummer-Tabellenname}		für Tabellen
%\label{eq:Kapitelnummer-Gleichungsname}	für Gleichungen
%\label{sec:Kapitelnummer}					für ganze Kapitel
%\cite{skript}
%
%Neues Programm: 
%
%Knight Rider (Leds leuchten hin und her)

%\lstinputlisting[caption=Beispiel 3, label = lst:bsp3, style=ST]{daten/CODE.txt}

%Ähnlich zu Listing \ref{lst:bsp1}, Listing \ref{lst:bsp2} und Listing \ref{lst:bsp3}


%\begin{figure}[H]
%	\centering
%	\def\svgwidth{250pt}
%	\input{images/Filename.png_tex}
%	\caption{Caption of a figure}
%	\label{fig:Lableofafigure}
%\end{figure}



